
\chapter{Glossar}

\begin{description}
  \item[\IX{Arteriosklerose}] Systemerkrankung der
  Arterien, welche zu Ablagerung von Fetten, Bindegewebe und Kalk in
  den Blutgefäßwänden und somit zur Schädigung und
  Veränderungen der Gefäßen führt.
  \item[\IX{bergen}]
  \item[\IX{Gestose}] Oberbegriff für durch die
  Schwangerschaft bedingte Erkrankungen. Man unterscheidet
  zwischen Gestosen der Frühschwangerschaft (im
  1. Trimenon, \I{Frühgestosen}) und der Spätschwangerschaft
  (3. Trimenon, \I{Spätgestosen}). Darunter fallen die
  \I{Hyperemesis gravidarum} (übermäßiges, anhaltendes
  Erbrechen), die EPH-Gestose bzw. Präekplampsie
  (\ref{topic:eph-gestose}), die Eklampsie
  (\ref{topic:eklampsie}) und das HELLP-Syndrom
  (\ref{topic:hellp-Syndrom}).
  \item[\IX{retten}]
  \item[\IX{Ruptur}] Riss, Einreissen, Durchbruch.
  \item[\IX{Sofortbehandlung}]
  \item[\IX{Sofortmaßnahmen, lebensrettende}]
\end{description}
